\documentclass[11pt]{amsart}
\usepackage[T1]{fontenc}
\usepackage{lmodern}
\usepackage{amsmath,amssymb,amsthm,mathtools}
\usepackage{xcolor}
\usepackage[colorlinks=true,linkcolor=blue!55!black,citecolor=blue!55!black,urlcolor=blue!55!black]{hyperref}

\newtheorem{theorem}{Theorem}[section]
\newtheorem{lemma}[theorem]{Lemma}
\newtheorem{proposition}[theorem]{Proposition}
\theoremstyle{definition}
\newtheorem{definition}[theorem]{Definition}
\newtheorem{remark}[theorem]{Remark}

\title[Odd-parameter periods for RESOCA 115]{Exact odd-parameter periods for reversible elementary second-order cellular automaton rule 115}
\author{Zijian Zeng}
\address{Institute of Computer Science and Digital Innovation, UCSI University, Kuala Lumpur 56000, Malaysia}
\email{zijianzeng@foxmail.com}
\email{1002266693@ucsiuniversity.edu.my}
\date{29 August 2026}
\hypersetup{
  pdftitle={Exact odd-parameter periods for reversible elementary second-order cellular automaton rule 115},
  pdfauthor={Zijian Zeng}
}

\begin{document}

\begin{abstract}
Formenti and Kamilya conjectured an exact period formula for the family
\(1{(01)}^k\) under reversible elementary second-order cellular automaton rule
115.  We prove their odd-parameter conjecture.  The proof encodes each pair of
consecutive time layers by four symbols, removes the uniform period-three
background, and conjugates the remaining dynamics to a phase-driven walk on
a finite properly three-edge-coloured tree.  A leaf-deletion argument shows
that this walk is one cycle, which supplies exact minimality rather than only
a return-time bound.  Consequently, for every \(i\geq0\),
\[
P(2i+1)=2^{4i+6}-1=4^{2i+3}-1.
\]
Independent exact computations verify the local conjugacy and small instances,
but the quantified conclusion follows from the symbolic proof.
\end{abstract}

\maketitle

\section{Introduction}

Formenti and Kamilya introduced and studied the reversible elementary
second-order cellular automaton induced by Wolfram rule~115
~\cite{FormentiKamilya2026}.  They proved that every finite initial pair is
periodic and investigated period functions for several structured families.
For the family whose current word is \(1{(01)}^k\) and whose memory is a single
central spike, their Conjecture~2 states that
\(P(2i+1)=4^{2i+3}-1\) for every \(i\geq0\).

We prove this conjecture by a nonlinear combinatorial conjugacy.  The natural
configuration coordinates do not give a linear map over \(\mathbb F_2\): the
local rule has algebraic normal form
\(1+x+\ell x+\ell xr\).  Instead, a phase-normalised two-layer symbol word
has the form \(3^{N-|v|}v\), where \(v\) is a reduced word over three colours.
The cellular update deletes or prepends one selected colour.  The allowed
words form a tree, and the three phases trace a single contour cycle around
that tree.

The source paper has a harmless endpoint typo in its generic word-placement
formula: a word of length \(|w|\) is indexed by
\(0,\ldots,|w|-1\), not through \(|w|\).  We make the centred initial
configuration explicit below.  We also use the standard least positive
return convention; without positivity, time zero would trivialise a period
definition.

\section{Rule and exact finite reduction}

All arithmetic in the cellular rule is over \(\mathbb F_2\).  Put
\[
\delta(\ell,x,r)=1+x+\ell x+\ell xr.
\]
Its Wolfram table is
\[
\begin{array}{c@{\quad}cccccccc}
(\ell x r)&111&110&101&100&011&010&001&000\\
\delta(\ell,x,r)&0&1&1&1&0&0&1&1,
\end{array}
\]
so this is rule 115.

For two configurations \(m,c\colon\mathbb Z\to\mathbb F_2\), define
\[
F(m,c)=(c,c'),\qquad
c'_j=m_j+\delta(c_{j-1},c_j,c_{j+1}).                 \tag{2.1}
\]
Let \(u_0=1\) and \(u_j=0\) for \(j\neq0\).  For \(k\geq0\), let
\[
c^{(k)}_j=
\begin{cases}
1,&-k\leq j\leq k\ \text{and}\ j+k\equiv0\pmod2,\\
0,&\text{otherwise}.
\end{cases}
\]
Thus sites \(-k,\ldots,k\) read as \(1{(01)}^k\).  The period is
\[
P(k)=\min\{p>0:F^p(u,c^{(k)})=(u,c^{(k)})\}.          \tag{2.2}
\]

Write the time layers as
\[
d^{-1}=u,\qquad d^0=c^{(k)},\qquad
d^{t+1}_j=d^{t-1}_j+
 \delta(d^t_{j-1},d^t_j,d^t_{j+1}).                  \tag{2.3}
\]
Let \(\beta_t=1\) for \(t\equiv1\pmod3\) and \(\beta_t=0\) otherwise,
also for \(t=-1\).  The uniform background pair cycles as
\[
(\beta_{t-1},\beta_t)=(0,0),(0,1),(1,0),(0,0),\ldots. \tag{2.4}
\]

\begin{lemma}[Exact exterior]\label{lem:exterior}
For \(k\geq1\) and every \(t\geq-1\),
\[
d^t_j=\beta_t\quad\text{if }j>k\text{ or }j\leq-k-2. \tag{2.5}
\]
Consequently a full return is equivalent to phase zero and return of both
layers on
\[
W_k=[-k-1,k].                                         \tag{2.6}
\]
\end{lemma}

\begin{proof}
Set \(a=-k\) and \(b=k\).  At the right receiver \(b+1\), the identities
\[
\delta(s,0,0)=1,\qquad \delta(s,1,1)=0
\]
are independent of the possibly non-background left neighbour \(s\).
Induction therefore proves (2.5) for \(j>b\).

For the left side put \(e_t=d^t_{a-1}\).  A simultaneous three-phase
induction gives
\[
e_{3q+1}=1,\qquad e_{3q}=e_{3q-1}\qquad(q\geq0),      \tag{2.7}
\]
starting from \(e_{-1}=e_0=0\).  At phase zero,
\(\delta(0,x,r)=1+x\) and the second identity in (2.7) give
\(e_{3q+1}=1\).  At phase two the same local identity gives
\(e_{3q+3}=e_{3q+2}\), which is the next instance of the second identity.

At site \(a-2\), the only phase at which the right neighbour can affect the
background update is phase one.  There the neighbourhood has centre and left
bits equal to one, while (2.7) gives the right bit one; hence
\(\delta(1,1,1)=0\), as required for the next background bit.  In the other
two phases, \(\delta(0,0,r)=1\), independently of \(r\).  This closes the
simultaneous induction at \(a-2\), and ordinary induction propagates the
background to every site further left.

Only the two background ghost cells \(a-2\) and \(b+1\) enter the update on
\(W_k\).  Conversely, at any site to the right of \(W_k\), the pair
\((\beta_{p-1},\beta_p)\) equals the initial pair \((0,0)\) exactly when
\(p\equiv0\pmod3\).  This proves the return equivalence.
\end{proof}

\section{A reduced-word tree}

Encode the two bits at site \(j\) and time \(t\) by
\[
q^t_j=2d^{t-1}_j+d^t_j\in Q=\{0,1,2,3\}.             \tag{3.1}
\]
Let \(\tau=(0\ 1\ 2)\) and fix \(\tau(3)=3\).  By (2.4), the exterior
symbol at phase \(p=t\bmod3\) is \(p\).  The site-symbol update is
\[
g(L,X,R)=2(X\bmod2)+
 \bigl(\lfloor X/2\rfloor+
 \delta(L\bmod2,X\bmod2,R\bmod2)\bigr)\bmod2.       \tag{3.2}
\]

Put \(N=2k+2=|W_k|\).  Let \(V_N\) contain the words
\(v=v_1\cdots v_m\) over \(\{0,1,2\}\) such that
\[
0\leq m\leq N,\qquad v_j\neq v_{j+1},\qquad
m>0\Rightarrow v_m\neq0,\qquad
m=N\Rightarrow v_1=0.                                \tag{3.3}
\]
The empty word is included.  Its normalised padded word is
\[
z(v)=3^{N-|v|}v.                                      \tag{3.4}
\]

For a colour \(a\in\{0,1,2\}\), define \(s_a\) by deleting the first
letter when it is \(a\), and otherwise attempting to prepend \(a\).  Perform
the operation if its result belongs to \(V_N\); if not, fix \(v\).  A valid
prepend is reversed by deletion, while an invalid prepend is a fixed point,
so \(s_a\) is an involution.

Join \(v\) to \(s_a(v)\) by an edge of colour \(a\) whenever the two words
differ.  Every nonempty word has the unique shorter neighbour obtained by
deleting its first letter, and repeated deletion reaches the empty word.
Thus the graph is connected.  On a hypothetical cycle, a word of maximal
length would need two shorter neighbours, contradicting uniqueness.  The graph
is therefore a tree, and at most one edge of any colour meets a vertex.

\begin{proposition}[Exact one-step conjugacy]\label{prop:conjugacy}
Let
\[
(a_0,a_1,a_2)=(1,0,2),                                \tag{3.5}
\]
and define
\[
\Psi(v,p)=\tau^p(z(v)).                                \tag{3.6}
\]
With exterior ghost symbol \(p\), one cellular step on \(W_k\) satisfies
\[
F_{W_k}(\Psi(v,p))=
\Psi(s_{a_p}(v),p+1\bmod3).                            \tag{3.7}
\]
\end{proposition}

\begin{proof}
For normalised symbols define
\[
h_p(L,X,R)=\tau^{-(p+1)}
g(\tau^{p}L,\tau^{p}X,\tau^{p}R).
\]
The normalised ghost is 0.  Direct substitution of rule 115 in (3.2) gives,
for every \((L,X,R)\in Q^3\),
\[
h_p(L,X,R)=
\begin{cases}
3,&X=a_p,\ L\in\{a_p,3\},\
     R\in\{0,1,2\}\setminus\{a_p\},\\
a_p,&X=3,\ L\in\{a_p,3\},\
     R\in\{0,1,2\}\setminus\{a_p\},\\
X,&\text{otherwise}.
\end{cases}                                           \tag{3.8}
\]
Thus only the boundary of the leading run of 3s can change, and the change is
\(a_p\leftrightarrow3\).

If the reduced word begins with \(a_p\), its first letter becomes 3 and is
deleted from the reduced word.  Otherwise the last leading 3 becomes \(a_p\)
and is prepended whenever the result satisfies (3.3).  Formula (3.8) also
handles every missing-edge stratum: an empty word accepts colours 1 and 2 but
not 0; a full word can only delete its forced first letter 0; and at length
\(N-1\), a prepend is possible exactly for colour 0.  These are precisely the
fixed cases in the definition of \(s_{a_p}\), proving (3.7).
\end{proof}

\section{The contour cycle and exact period}

\begin{lemma}[Three-colour contour]\label{lem:contour}
Let \(T\) be a finite tree with a proper edge colouring by three colours.  Let
\(s_a\) cross the incident edge of colour \(a\), or fix the vertex when that
edge is missing.  If \((a_0,a_1,a_2)\) is any ordering of the colours, then
\[
\Phi(v,p)=(s_{a_p}(v),p+1\bmod3)                      \tag{4.1}
\]
is one cycle on \(V(T)\times\mathbb Z/3\mathbb Z\).
\end{lemma}

\begin{proof}
Induct on \(|V(T)|\).  For one vertex, (4.1) is its phase 3-cycle.  Remove a
leaf \(u\), attached to \(v\) by the edge of colour \(a_p\).  In the smaller
tree the missing-edge transition is
\((v,p)\to(v,p+1)\).  Restoring \(u\) replaces exactly this transition by
\[
(v,p)\to(u,p+1)\to(u,p+2)\to(u,p)\to(v,p+1).          \tag{4.2}
\]
The two intermediate selected colours differ from \(a_p\), so the
corresponding moves fix the leaf.  Hence the three new phase states are
spliced into the unique old cycle, and the result is again one cycle.
\end{proof}

It remains to count the tree and identify the initial state.  For
\(1\leq m<N\), there are \(2^m\) words of length \(m\) satisfying (3.3):
choose the final nonzero letter in two ways and each preceding unequal letter
in two ways.  For length \(N\), the first letter is 0.  The number ending
nonzero is
\[
A_N=\frac{2^N+2{(-1)}^N}{3},                           \tag{4.3}
\]
as follows from the two-state recurrence according to whether the last letter
is zero or nonzero.  Since \(N=2k+2\) is even,
\[
|V_N|=1+\sum_{m=1}^{N-1}2^m+A_N
=\frac{2^{N+2}-1}{3}.                                 \tag{4.4}
\]

\begin{theorem}\label{thm:main}
For every integer \(i\geq0\), the least positive full two-layer return time of
the memory \(u\) and current word \(1{(01)}^{2i+1}\) under RESOCA rule 115 is
\[
P(2i+1)=2^{4i+6}-1=4^{2i+3}-1.
\]
\end{theorem}

\begin{proof}
Let \(k=2i+1\).  At phase zero the symbol word on \([-k-1,k]\) begins with
0 and ends with 1.  Away from the origin it follows the alternating current
word.  Since \(k\) is odd, the current bit at site zero is 0, while the memory
spike changes its symbol to 2.  Thus adjacent symbols remain unequal, and the
entire length-\(N\) word is a vertex of \(V_N\).

By Proposition~\ref{prop:conjugacy} and
Lemma~\ref{lem:contour}, every phase-tagged tree state has least period
\(3|V_N|\).  The physical encoding is injective: the exterior background
recovers the phase, inverse rotation recovers \(z(v)\), and deleting its
leading run of 3s recovers \(v\).  Lemma~\ref{lem:exterior} therefore makes
the tree period equal to the full-lattice period.  Using (4.4),
\[
P(k)=3|V_N|=2^{N+2}-1=2^{2k+4}-1.
\]
Substitution of \(k=2i+1\) gives the formula.  The one-cycle conclusion also
shows directly that no smaller positive time returns.
\end{proof}

\section{Independent computational verification}

The proof above is all-parameter and does not infer a theorem from finite
data.  Exact computations were nevertheless used to attack its fragile local
and boundary claims.  One implementation reconstructs (3.2) from the integer
115 and checks every vertex and phase through \(N=12\).  A separately written
implementation uses a literal truth table and exhausts all
\(3\cdot4^3=192\) normalised local contexts and every tree state through
\(N=14\).  At \(N=14\) it finds 21,845 vertices and one cycle on 65,535 phase
states.

Two further simulators enumerate positive times until the first full return on
the exact window.  One uses dense tuples and the truth table, while the other
uses bit-packed layers and the algebraic normal form.  Both give
\[
P(1)=63,\quad P(3)=1023,\quad P(5)=16383,\quad
P(7)=262143.
\]
The serialized certificate validates the formal statement and fixed
implementation files, rejects malformed input, and reruns these verification
layers with one command.  These computations certify the displayed finite
checks, not the inductive passage to arbitrary \(i\).

\section*{Data and code availability}

The accompanying research archive contains the formal statement, human proof,
exact certificate, independent verifier, and deterministic tests.  The
principal verification command is
\begin{center}
\texttt{python3 src/verify\_tree\_certificate.py}.
\end{center}
No floating-point arithmetic, randomness, or proof assistant is used.

\bibliographystyle{amsplain}
\bibliography{references}

\end{document}
